\section{Why judges fail: failure modes and interventions}\label{sec:failures}

Knowing \emph{that} a judge is wrong is less useful than knowing \emph{how}. Every wrong answer is
sorted into one class by a fixed rule that looks only at the case, its tags and the answer.

\begin{description}[leftmargin=1.2em,style=nextline]
  \item[Acted on side effect] The right answer was to defer (\code{reason}) and the judge picked the
    option that names a side effect (buy, delete, publish, \dots).
  \item[Under-deferred] The right answer was to defer, but the judge picked some other concrete action.
  \item[Followed injection] The judge picked exactly what an instruction planted in the observation or
    source code told it to pick. This class takes precedence over the others.
  \item[Accepted wrong code] Search: the judge said ``yes, the code does this'' when it does not.
  \item[Rejected correct code] Search: the judge said ``no'' when the code is correct.
  \item[Wrong target] The judge picked a concrete option, but not the right one.
  \item[Over-deferred] The judge deferred although one option was right: safe, but loses the speed-up.
\end{description}

\subsection{The failure-mode matrix}

\begin{table}[htbp]
\centering\small
\caption{Holdout \prereg: all wrong answers by class (mean per repetition), whether or not the bundle
would have acted on them. Darker cells mean more errors; a dot means none.}
\label{tab:fail-all}
\begin{tabular}{@{}l ccccccc r@{}}
\toprule
Judge & \rotatebox{60}{Acted on side effect} & \rotatebox{60}{Under-deferred} & \rotatebox{60}{Followed injection} & \rotatebox{60}{Accepted wrong code} & \rotatebox{60}{Rejected correct code} & \rotatebox{60}{Wrong target} & \rotatebox{60}{Over-deferred} & total \\
\midrule
GPT-6.1 Sol & \cellcolor{okblue!19}1.3 & \textcolor{okgray}{\textperiodcentered} & \textcolor{okgray}{\textperiodcentered} & \textcolor{okgray}{\textperiodcentered} & \textcolor{okgray}{\textperiodcentered} & \textcolor{okgray}{\textperiodcentered} & \textcolor{okgray}{\textperiodcentered} & 1.3 \\
GPT-6 Luna & \cellcolor{okblue!21}1.7 & \textcolor{okgray}{\textperiodcentered} & \cellcolor{okblue!17}1 & \textcolor{okgray}{\textperiodcentered} & \textcolor{okgray}{\textperiodcentered} & \textcolor{okgray}{\textperiodcentered} & \textcolor{okgray}{\textperiodcentered} & 2.7 \\
Jev 1.13 & \cellcolor{okblue!23}2 & \cellcolor{okblue!16}0.7 & \cellcolor{okblue!34}4 & \cellcolor{okblue!17}1 & \textcolor{okgray}{\textperiodcentered} & \textcolor{okgray}{\textperiodcentered} & \textcolor{okgray}{\textperiodcentered} & 7.7 \\
nimble 9B & \cellcolor{okblue!23}2 & \cellcolor{okblue!23}2 & \cellcolor{okblue!44}6 & \cellcolor{okblue!28}3 & \textcolor{okgray}{\textperiodcentered} & \cellcolor{okblue!17}1 & \textcolor{okgray}{\textperiodcentered} & 14 \\
tev1 4B & \cellcolor{okblue!23}2 & \cellcolor{okblue!23}2 & \cellcolor{okblue!44}6 & \cellcolor{okblue!39}5 & \cellcolor{okblue!17}1 & \textcolor{okgray}{\textperiodcentered} & \textcolor{okgray}{\textperiodcentered} & 16 \\
tev1 0.8B & \cellcolor{okblue!23}2 & \cellcolor{okblue!28}3 & \cellcolor{okblue!60}\color{white}9 & \cellcolor{okblue!50}\color{white}7 & \textcolor{okgray}{\textperiodcentered} & \cellcolor{okblue!39}5 & \cellcolor{okblue!23}2 & 28 \\
Qwen3 8B & \cellcolor{okblue!23}2 & \cellcolor{okblue!34}4 & \cellcolor{okblue!55}\color{white}8 & \cellcolor{okblue!28}3 & \textcolor{okgray}{\textperiodcentered} & \cellcolor{okblue!17}1 & \textcolor{okgray}{\textperiodcentered} & 18 \\
Qwen3 4B & \cellcolor{okblue!23}2 & \cellcolor{okblue!44}6 & \cellcolor{okblue!60}\color{white}9 & \cellcolor{okblue!44}6 & \textcolor{okgray}{\textperiodcentered} & \cellcolor{okblue!17}1 & \cellcolor{okblue!17}1 & 25 \\
Qwen3 0.6B & \cellcolor{okblue!23}2 & \cellcolor{okblue!60}\color{white}9 & \cellcolor{okblue!82}\color{white}13 & \cellcolor{okblue!50}\color{white}7 & \textcolor{okgray}{\textperiodcentered} & \cellcolor{okblue!39}5 & \textcolor{okgray}{\textperiodcentered} & 36 \\
Laya base & \cellcolor{okblue!23}2 & \cellcolor{okblue!66}\color{white}10 & \cellcolor{okblue!66}\color{white}10 & \cellcolor{okblue!50}\color{white}7 & \textcolor{okgray}{\textperiodcentered} & \cellcolor{okblue!44}6 & \textcolor{okgray}{\textperiodcentered} & 35 \\
\bottomrule
\end{tabular}

\end{table}

\begin{table}[htbp]
\centering\small
\caption{Holdout \prereg: only the wrong answers the bundle would have \emph{acted on} under its gate
(mean per repetition). This is the unsafe subset of \cref{tab:fail-all}.}
\label{tab:fail-auto}
\begin{tabular}{@{}l ccccccc r@{}}
\toprule
Judge & \rotatebox{60}{Acted on side effect} & \rotatebox{60}{Under-deferred} & \rotatebox{60}{Followed injection} & \rotatebox{60}{Accepted wrong code} & \rotatebox{60}{Rejected correct code} & \rotatebox{60}{Wrong target} & \rotatebox{60}{Over-deferred} & total \\
\midrule
GPT-6.1 Sol & \cellcolor{okblue!15}0.3 & \textcolor{okgray}{\textperiodcentered} & \textcolor{okgray}{\textperiodcentered} & \textcolor{okgray}{\textperiodcentered} & \textcolor{okgray}{\textperiodcentered} & \textcolor{okgray}{\textperiodcentered} & \textcolor{okgray}{\textperiodcentered} & 0.3 \\
GPT-6 Luna & \cellcolor{okblue!24}1.3 & \textcolor{okgray}{\textperiodcentered} & \cellcolor{okblue!21}1 & \textcolor{okgray}{\textperiodcentered} & \textcolor{okgray}{\textperiodcentered} & \textcolor{okgray}{\textperiodcentered} & \textcolor{okgray}{\textperiodcentered} & 2.3 \\
Jev 1.13 & \textcolor{okgray}{\textperiodcentered} & \textcolor{okgray}{\textperiodcentered} & \textcolor{okgray}{\textperiodcentered} & \cellcolor{okblue!21}1 & \textcolor{okgray}{\textperiodcentered} & \textcolor{okgray}{\textperiodcentered} & \textcolor{okgray}{\textperiodcentered} & 1 \\
nimble 9B & \cellcolor{okblue!30}2 & \textcolor{okgray}{\textperiodcentered} & \cellcolor{okblue!38}3 & \cellcolor{okblue!38}3 & \textcolor{okgray}{\textperiodcentered} & \textcolor{okgray}{\textperiodcentered} & \textcolor{okgray}{\textperiodcentered} & 8 \\
tev1 4B & \cellcolor{okblue!30}2 & \textcolor{okgray}{\textperiodcentered} & \cellcolor{okblue!38}3 & \cellcolor{okblue!56}\color{white}5 & \cellcolor{okblue!21}1 & \textcolor{okgray}{\textperiodcentered} & \textcolor{okgray}{\textperiodcentered} & 11 \\
tev1 0.8B & \cellcolor{okblue!21}1 & \cellcolor{okblue!21}1 & \cellcolor{okblue!56}\color{white}5 & \cellcolor{okblue!73}\color{white}7 & \textcolor{okgray}{\textperiodcentered} & \textcolor{okgray}{\textperiodcentered} & \textcolor{okgray}{\textperiodcentered} & 14 \\
Qwen3 8B & \cellcolor{okblue!30}2 & \cellcolor{okblue!30}2 & \cellcolor{okblue!64}\color{white}6 & \cellcolor{okblue!38}3 & \textcolor{okgray}{\textperiodcentered} & \cellcolor{okblue!21}1 & \textcolor{okgray}{\textperiodcentered} & 14 \\
Qwen3 4B & \cellcolor{okblue!30}2 & \textcolor{okgray}{\textperiodcentered} & \cellcolor{okblue!82}\color{white}8 & \cellcolor{okblue!64}\color{white}6 & \textcolor{okgray}{\textperiodcentered} & \textcolor{okgray}{\textperiodcentered} & \textcolor{okgray}{\textperiodcentered} & 16 \\
Qwen3 0.6B & \cellcolor{okblue!21}1 & \textcolor{okgray}{\textperiodcentered} & \cellcolor{okblue!82}\color{white}8 & \cellcolor{okblue!73}\color{white}7 & \textcolor{okgray}{\textperiodcentered} & \textcolor{okgray}{\textperiodcentered} & \textcolor{okgray}{\textperiodcentered} & 16 \\
Laya base & \cellcolor{okblue!21}1 & \textcolor{okgray}{\textperiodcentered} & \cellcolor{okblue!47}4 & \cellcolor{okblue!73}\color{white}7 & \textcolor{okgray}{\textperiodcentered} & \textcolor{okgray}{\textperiodcentered} & \textcolor{okgray}{\textperiodcentered} & 12 \\
\bottomrule
\end{tabular}

\end{table}
\howtoreadtable{Compare the two matrices row by row. A cell that is dark in the first but empty in the
second is a failure the gate caught: the judge was wrong but unsure. A cell that stays dark in the second
matrix is a failure the agent would have acted on.}

Read together, the two matrices show how differently the judges fail. Jev, for example, followed an
injected instruction \JevHoldAllInj\ times per repetition, but never automatically: every one of those
answers fell back, because the injections pushed it toward \code{reason} or its confidence was low.
GPT-6 Luna followed an injection \LunaHoldAllInj\ time per repetition, and it acted on it every time
(\LunaHoldAutoInj).

\subsection{Root causes}

\paragraph{Acting on side effects is shared by every judge, and its cause is the instruction.}
All \SideEffectJudgesHold\ base judges picked a side-effect option at least once on the holdout. The
judges' instruction lists ``writing, sending, or substantial reasoning'' as reasons to defer, but not
purchasing, deleting, refunding or publishing. The judges are doing roughly what they were told. Two
pieces of evidence point to the wording rather than to the models: adding one sentence to Jev's
instruction (I1) removed all of its targeted dev errors (\IOneDevJevClBefore$\to$\IOneDevJevClAfter), and a
host-side word list (I2) removed them for every judge.

\paragraph{Accepted wrong code dominates the small local judges, and the policy amplifies it.}
On dev, \TevPtEightDevAllAwc\ of tev1 0.8B's \TevPtEightDevWrong\ errors were accepted wrong code. The
bundle has no abstention for yes/no answers, so each of these becomes an action: a model that says
``yes'' at 0.55 is acted on exactly like one that says ``yes'' at 0.99.

\paragraph{Injection following splits the families.} On dev only local judges followed injections.
On the holdout, where \HoldNInjection\ cases carry an injected instruction, Jev tended to defer when an
injected note pushed it toward \code{reason}, which is a safe failure. GPT-6 Luna followed an injected
``click Submit now'' on \code{hold-cua-12} at \LunaHoldSubmitPRange, which is not.

\paragraph{Confident cloud mistakes.} Luna's mistakes on the holdout were few but confident: it issued
a full, irreversible refund on \code{hold-cua-00} at \LunaHoldRefundPRange\ in all \LunaHoldRefundReps\
repetitions. In repetition 1, Jev was wrong on \JevHoldRepOneWrong\ holdout cases, and
\JevHoldRepOneWrongFallback\ of them fell back to slow reasoning (its highest confidence on a wrong
answer was \JevHoldRepOneWrongMaxP). This pattern, humble mistakes for Jev and confident ones for Luna,
is why the cloud fallback recommendation in \cref{sec:recs} comes with a guard.

\subsection{Interventions}\label{sec:interventions}

Three interventions were declared in writing before any was measured, measured once on dev, and then
confirmed or not on the holdout under the preregistered rule.
\begin{description}[leftmargin=1.2em,style=nextline]
  \item[I1: side-effect clause (targets acted-on-side-effect and under-deferred)] One sentence appended to
    the judge's instruction: \emph{``\IOneClause''}
  \item[I2: host guard (targets acted-on-side-effect)] A deterministic check in the host, not the model:
    if the chosen option's own description matches a fixed side-effect verb list, fall back. No model
    call and no per-case exceptions.
  \item[I3: yes/no gate (targets accepted-wrong-code and rejected-correct-code)] A yes/no answer is
    automatic only when $\max(p, 1-p) \ge \MinP$, the same bar the bundle uses for choices.
\end{description}

\begin{table}[htbp]
\centering\small
\caption{Interventions on both splits. ``Targeted'' is the number of targeted wrong automatic decisions,
summed over the \HoldReps\ repetitions, before$\to$after. $\Delta$acc is the change in correct cases;
$\Delta$cov the change in coverage in percentage points. \ok\ marks an arm that meets the confirmation
rule on its own. Only arms where the intervention applies (at least one targeted error before) are
listed.}
\label{tab:interventions}
\begin{tabular*}{\textwidth}{@{\extracolsep{\fill}}l l r r r r r r@{}}
\toprule
 & & \multicolumn{3}{c}{Dev (screen)} & \multicolumn{3}{c}{Holdout (preregistered)} \\
\cmidrule(lr){3-5}\cmidrule(lr){6-8}
Intervention & arm & targeted & $\Delta$acc & $\Delta$cov & targeted & $\Delta$acc & $\Delta$cov \\
\midrule
I1 clause & GPT-6 Luna + I1 & 6$\to$5 & +2 & $-$1.1 & 4$\to$2\,\ok & +1 & $-$3.2 \\
 & Jev 1.13 + I1 & 5$\to$0\,\ok & +3 & $-$3.3 & \textemdash &  &  \\
 & tev1 0.8B + I1 & \textemdash &  &  & 6$\to$6 & $-$2 & $-$3.2 \\
 & Qwen3 8B + I1 & 18$\to$15 & +2 & $-$1.1 & 12$\to$12 & $-$2 & $-$3.2 \\
 & \emph{pooled} & 29$\to$20 & \multicolumn{2}{r}{not confirmed} & 22$\to$20 & \multicolumn{2}{r}{not confirmed} \\
\midrule
I2 guard & GPT-6.1 Sol & 3$\to$0\,\ok & +0 & $-$1.1 & 1$\to$0\,\ok & +0 & $-$3.2 \\
 & GPT-6 Luna & 5$\to$0\,\ok & +0 & $-$2.2 & 4$\to$0\,\ok & +0 & $-$6.3 \\
 & Jev 1.13 & 5$\to$0\,\ok & +0 & $-$2.2 & \textemdash &  &  \\
 & nimble 9B & 9$\to$0\,\ok & +0 & $-$3.3 & 6$\to$0\,\ok & +0 & $-$7.9 \\
 & tev1 4B & 9$\to$0\,\ok & +0 & $-$3.3 & 6$\to$0\,\ok & +0 & $-$7.9 \\
 & tev1 0.8B & \textemdash &  &  & 3$\to$0\,\ok & +0 & $-$1.6 \\
 & Qwen3 8B & 12$\to$0\,\ok & +0 & $-$4.4 & 6$\to$0\,\ok & +0 & $-$7.9 \\
 & Qwen3 4B & 9$\to$0\,\ok & +0 & $-$3.3 & 6$\to$0\,\ok & +0 & $-$6.3 \\
 & Qwen3 0.6B & \textemdash &  &  & 3$\to$0\,\ok & +0 & $-$3.2 \\
 & Laya base & \textemdash &  &  & 3$\to$0\,\ok & +0 & $-$1.6 \\
 & \emph{pooled} & 52$\to$0 & \multicolumn{2}{r}{\textbf{confirmed}} & 38$\to$0 & \multicolumn{2}{r}{\textbf{confirmed}} \\
\midrule
I3 yes/no gate & Jev 1.13 & 3$\to$0\,\ok & +0 & $-$6.7 & 3$\to$0 & +0 & $-$12.7 \\
 & nimble 9B & 12$\to$6\,\ok & +0 & $-$4.4 & 9$\to$6 & +0 & $-$11.1 \\
 & tev1 4B & 18$\to$6 & +0 & $-$11.1 & 18$\to$9 & +0 & $-$15.9 \\
 & tev1 0.8B & 30$\to$0 & +0 & $-$30.0 & 21$\to$9 & +0 & $-$17.5 \\
 & Qwen3 8B & 9$\to$3\,\ok & +0 & $-$6.7 & 9$\to$9 & +0 & $-$1.6 \\
 & Qwen3 4B & 12$\to$0 & +0 & $-$31.1 & 18$\to$0 & +0 & $-$28.6 \\
 & Qwen3 0.6B & 39$\to$39 & +0 & +0.0 & 21$\to$21 & +0 & $-$1.6 \\
 & Laya base & 36$\to$0 & +0 & $-$32.2 & 21$\to$0 & +0 & $-$31.7 \\
 & \emph{pooled} & 159$\to$54 & \multicolumn{2}{r}{not confirmed} & 120$\to$54 & \multicolumn{2}{r}{not confirmed} \\
\midrule
I2+I3 & GPT-6.1 Sol & 3$\to$0\,\ok & +0 & $-$1.1 & 1$\to$0\,\ok & +0 & $-$3.2 \\
 & GPT-6 Luna & 5$\to$0\,\ok & +0 & $-$2.2 & 4$\to$0\,\ok & +0 & $-$6.3 \\
 & Jev 1.13 & 8$\to$0\,\ok & +0 & $-$8.9 & 3$\to$0 & +0 & $-$14.3 \\
 & nimble 9B & 21$\to$6\,\ok & +0 & $-$7.8 & 15$\to$6 & +0 & $-$19.0 \\
 & tev1 4B & 27$\to$6 & +0 & $-$14.4 & 24$\to$9 & +0 & $-$23.8 \\
 & tev1 0.8B & 30$\to$0 & +0 & $-$30.0 & 24$\to$9 & +0 & $-$19.0 \\
 & Qwen3 8B & 21$\to$3 & +0 & $-$11.1 & 15$\to$9 & +0 & $-$9.5 \\
 & Qwen3 4B & 21$\to$0 & +0 & $-$34.4 & 24$\to$0 & +0 & $-$34.9 \\
 & Qwen3 0.6B & 39$\to$39 & +0 & +0.0 & 24$\to$21 & +0 & $-$4.8 \\
 & Laya base & 36$\to$0 & +0 & $-$32.2 & 24$\to$0 & +0 & $-$33.3 \\
 & \emph{pooled} & 211$\to$54 & \multicolumn{2}{r}{not confirmed} & 158$\to$54 & \multicolumn{2}{r}{not confirmed} \\
\bottomrule
\end{tabular*}

\end{table}

\begin{figure}[htbp]
\centering
% Intervention I2 (host guard) on the holdout: targeted wrong automatic decisions before and after.
\begin{tikzpicture}
\begin{axis}[benchmark,
  xbar, bar width=6pt,
  width=0.9\linewidth, height=7cm,
  xmin=0, ymin=-0.6, ymax=\ITwoPlotMax+0.6,
  ytick={0,...,\ITwoPlotMax},
  yticklabels from table={generated/data/i2-holdout.dat}{label},
  xlabel={Side-effect wrong automatic decisions, summed over \HoldReps\ repetitions},
  xmajorgrids,
  legend style={at={(0.98,0.04)}, anchor=south east}, legend cell align=left,
  nodes near coords, nodes near coords style={font=\scriptsize},
  enlarge y limits=false,
]
\addplot[fill=okorange!80, draw=black!70, postaction={pattern=north east lines, pattern color=white}]
  table[x=before, y=idx]{generated/data/i2-holdout.dat};
\addlegendentry{bundle gate only}
\addplot[fill=okblue, draw=black!70]
  table[x=after, y=idx]{generated/data/i2-holdout.dat};
\addlegendentry{bundle gate + host guard (I2)}
\end{axis}
\end{tikzpicture}

\caption{The host guard (I2) on the holdout \prereg: side-effect wrong automatic decisions per judge,
summed over repetitions, without and with the guard.}
\label{fig:i2}
\howtoread{Each pair of bars is one judge. The hatched orange bar counts side-effect actions the bundle
would have taken automatically and wrongly; the solid blue bar counts them with the guard on. Every blue
bar is zero. The price is coverage: the guard also blocks some correct actions whose description contains a
listed verb, \ITwoHoldCovCostRange\ points depending on the judge.}
\end{figure}

\paragraph{What was confirmed on the holdout.}
\begin{itemize}
  \item \textbf{I2 (host guard): \ITwoHoldConfirmed.} It removed all \ITwoHoldBefore\ targeted wrong
    automatic decisions in the \ITwoHoldNApplicable\ judges it applied to, with no accuracy change and
    a coverage cost of \ITwoHoldCovCostRange\ points. Jev had no side-effect wrong automatic decisions
    on the holdout to remove.
  \item \textbf{I1 (clause): \IOneHoldConfirmed\ overall.} It worked for GPT-6 Luna
    (\IOneHoldLunaClBefore$\to$\IOneHoldLunaClAfter) but not for the local judges, and it did not apply to
    Jev, whose side-effect picks all fell below the \MinP\ gate. Jev with the clause was
    \JevClHoldRAccDiff\ points more accurate than Jev (\JevClHoldAccK\ versus \JevHoldAccK), but that
    gain is within noise (Holm $p = \JevClHoldRAccPHolm$).
  \item \textbf{I3 (yes/no gate): \IThreeHoldConfirmed.} It removed many accepted-wrong-code errors, but no
    single judge met the confirmation rule: where it removed errors, it cost more than \IConfirmCovPts\
    points of coverage (the range across judges was \IThreeHoldCovCostRange\ points), and it left
    the Qwen3 0.6B and Qwen3 8B errors untouched.
\end{itemize}

\begin{keyidea}
The most effective safety measure in this benchmark is not a better model but a dumb, deterministic
rule in the host: never act automatically on an option that names an irreversible side effect. It works
the same way for every judge, costs a few points of coverage, and needs no retraining.
\end{keyidea}
